\documentclass[12pt,letterpaper]{article}
\usepackage[T1]{fontenc}
\usepackage{mathptmx}
\usepackage{courier}
\usepackage[letterpaper,margin=1in]{geometry}
\usepackage{amsmath,amssymb}
\usepackage{microtype}
\usepackage{booktabs}
\usepackage{array}
\usepackage{enumitem}
\usepackage[hidelinks]{hyperref}

\newcommand{\F}{\mathbb{F}}
\newcommand{\Z}{\mathbb{Z}}
\DeclareMathOperator{\PGL}{PGL}
\DeclareMathOperator{\PSL}{PSL}
\DeclareMathOperator{\GL}{GL}
\DeclareMathOperator{\Ind}{Ind}

\title{Referee report on the hand-edited Draft 1 of\\
\emph{The representation theory of finite commutation transitive groups}\\
(file \texttt{commutation-transitive-groups2.tex})}
\author{Referee: Claude (Anthropic, Claude Opus 5.5)}
\date{November 17, 2026}

\begin{document}
\maketitle

\section*{What was read}

I read the whole uploaded file (1,810 lines) and compiled it with pdflatex.
For comparison I used Draft~1 as it stands in the repository (commit
\texttt{bd9e320}) and the house style sheet
\texttt{joyner-math-paper-style-guide.md} (Draft~2). I did not consult the
full texts of the cited papers. The revision log of Draft~2 lists the
sources consulted for the bibliography.

Line numbers below refer to the uploaded file. Each item has a verdict
(correct, misprint, fixable gap or false) and a verification level. Level~1
is an independent derivation. Level~2 is a check against a source I
consulted. Level~3 is reliance on memory. Every item has been dealt with in
Draft~2, as recorded in its revision log.

\section*{Summary}

The hand edits remove most comparisons with the 2000 manuscript and recast
several results as answers to natural conjectures. That is a good
expository device. The edits left three broken cross-references. They also
left one false statement in the introduction (some of the ``false''
conjectures are true), and one incorrect justification, in
Question~10.5. The mathematics of the theorems and proofs is unchanged from
Draft~1, and I found no new errors in it. The remaining points concern
precision and style.

\section{Correctness}

\begin{enumerate}[label=\textbf{C\arabic*.},leftmargin=*,itemsep=4pt]
\item \textbf{Broken references.} The introduction (line 144) cites
  Counterexamples \verb|cex:formula| and \verb|cex:supports|. The summary
  table (lines 1569, 1574, 1582) also cites \verb|rem:gap|. All three items
  were deleted, so pdflatex prints ``??''. \emph{Verdict:} misprint.
  \emph{Level:} 1 (compiled).

\item \textbf{Question 10.5, the profinite question (lines 1623--1630).} It
  says that a finite quotient of a CT group need not be CT, ``compare
  Counterexample~5.18''. That counterexample is about commutation transitive
  groups and the center, not about CT groups. Moreover every quotient of a
  \emph{finite} CT group is CT, by the result of Wu cited in the same draft.
  The claim itself is true for infinite CT groups, for the right reason.
  Free groups are CT, and every finite group is a quotient of a free group.
  \emph{Verdict:} false justification of a true statement. \emph{Level:} 1.

\item \textbf{The introduction misdescribes the conjectures (lines
  140--144).} It says the conjectural statements are false and each is
  refuted by a counterexample. But the conjecture before
  Proposition~7.4 (dimension at most~$4$) is true, and
  Proposition~7.4 proves it with the sharp bound~$3$. The ``classification''
  conjecture in Section~8 is also true (Theorem~8.3). The conjecture on
  $\PSL(2,\F_q)$ is refuted by a proposition, not a counterexample.
  \emph{Verdict:} false as written. \emph{Level:} 1.

\item \textbf{The third conjecture of Section 5.5 (lines 932--934)}, that
  ``commutation in $G/Z$ is `the same as' commutation in $G$'', is not a
  mathematical statement. The counterexample that follows refutes a precise
  statement: if $G$ is commutation transitive, then so is $G/Z(G)$.
  \emph{Verdict:} fixable gap. \emph{Level:} 1.

\item \textbf{The conjecture of Section 8 (lines 1244--1246)} uses the
  symbol $G^*$, which the draft no longer defines. As stated (all
  representations are linear combinations of monomial representations), it
  holds for \emph{every} finite group, by Brauer's induction theorem. So it
  says nothing specific about CT groups. The content of Theorem~8.3 is that
  monomial characters induced from \emph{centralizers} suffice. The
  conjecture should say this. \emph{Verdict:} fixable gap. \emph{Level:} 1
  for the logic, 3 for the location of Brauer's theorem in Serre's book.

\item \textbf{Wu's quotient-closure statement (Counterexample 5.18, lines
  944--946).} The statement rests on secondary summaries of Wu's paper, not
  on the paper itself. Its connection to the counterexample is also unclear.
  For CT groups the analogue about $G/Z$ says nothing, because a non-abelian
  CT group has trivial center. The meaningful statement is that every
  quotient of a locally finite CT group is CT. It belongs in a
  ``Remark~2026.$n$'' (style sheet, rule~6.6). \emph{Verdict:} fixable gap.
  \emph{Level:} 2 (secondary sources only).

\item \textbf{Title of Subsection 5.3, ``Centralizers are malnormal''
  (line 708).} Lemma~5.8(c) shows that two conjugate centralizers are equal
  or meet in~$Z$. A centralizer is malnormal only when $Z=1$ and the
  centralizer is self-normalizing. \emph{Verdict:} misprint (inaccurate
  title). \emph{Level:} 1.

\item \textbf{Bibliography and literature review.}
  \begin{itemize}[itemsep=1pt]
  \item The ``commutative transitive kernel'' paper is cited as ``cited in
    2000 as `to appear in Algebra Colloquium'\,''. It appeared in
    \emph{Algebra Colloq.}\ \textbf{4} (1997), 141--152. So it belongs in
    the pre-2000 part of the literature review. \emph{Level:} 2.
  \item Fine--Rosenberger 2009 lacks its page range, 112--130.
    \emph{Level:} 2.
  \item The entry \texttt{Jo00} is no longer cited. \emph{Level:} 1.
  \item ``The first serious attack on Burnside's conjecture'' (line 179)
    overstates the claim. ``An early step toward'' is safe. \emph{Level:} 3.
  \end{itemize}
  \emph{Verdict:} misprints.
\end{enumerate}

\paragraph{Checked and correct.} I re-read every proof in the draft:
Lemmas~3.1, 3.6, 3.9, 3.13, 5.8, 8.2 and~9.4, Propositions~5.2, 5.5, 5.6,
5.9, 5.13, 5.14, 5.17, 6.1, 6.6, 7.3, 7.4, 9.2 and~9.7,
Theorems~6.3, 8.3 and~9.5, and Corollaries~6.4 and~9.8. All are correct
(level~1). The numbers are those of the uploaded file. I reran the
verification script \texttt{comm\_tr\_checks.py}. Its output is identical to
the output for Draft~1, and agrees with every computational claim in the
draft (level~1).

\section{Readability}

\begin{enumerate}[label=\textbf{R\arabic*.},leftmargin=*,itemsep=4pt]
\item \textbf{The conjectures are informal and unnumbered.} They appear as
  ``here's a natural \emph{Conjecture}'', ``Is the following classification
  true?'' and so on (lines 873, 891, 932, 1125, 1212, 1244), and cannot be
  cross-referenced. One counterexample starts ``You can't transport\dots''
  (line 938). I suggest a numbered environment, ``Naive conjecture'', with a
  precise statement and a one-sentence motivation for each.
\item \textbf{Lines 213--214.} ``draafted'' is a typo. The paragraph is a
  single sentence, and it fits better in the opening of the literature
  review, where it explains the split at 2000.
\item \textbf{Line 783.} ``Both directions of `if' can be checked directly''
  is unclear. The intended meaning is that the ``if'' direction of
  Theorem~5.11 can be checked directly for cases (ii) and~(iii).
\end{enumerate}

\section{Style sheet}

\begin{enumerate}[label=\textbf{S\arabic*.},leftmargin=*,itemsep=4pt]
\item \textbf{Rule 6.2.} The subtitle still reads ``Draft 1''.
\item \textbf{Rule 6.4 (revision history out of the body).} Comparisons with
  the 2000 manuscript remain at lines 127 (``studied in 2000''), 297--298
  (left cosets), 878 (Counterexample~5.16), 1552 and~1558 (summary table),
  1691 and~1694 (Acknowledgments) and~1737 (bibliography). The title-page
  footnote, the date line and the literature-review split at 2000 should
  stay, since rules~6.2, 6.3 and~2.3 require them.
\item \textbf{Rule 3.2 (orienting paragraphs).} The following begin without
  one: Subsections~3.1 and~3.3, Section~4 (Notation), Subsections~5.3
  and~5.4, Subsections~6.1 and~6.2, Subsection~9.1, and Section~11
  (Glossary).
\end{enumerate}

\section*{Recommendation}

All points are local and have been fixed in Draft~2. After these
corrections I see no remaining mathematical errors. The ``naive
conjecture'' framing makes the note easier to read than Draft~1.

\end{document}
